% ============================================================
% Projective Dynamic Logo (PDL) --- Document D61
% Gauge Dynamics as a C4 Survival Theorem:
% Derivation of the Minimal Covariant Derivative from C1--C4
% C. Laubscher, 2026
% ============================================================
\documentclass[11pt,a4paper]{article}

% ---- Encoding and fonts ------------------------------------
\usepackage[utf8]{inputenc}
\usepackage[T1]{fontenc}
\usepackage{lmodern}
\usepackage{newtxtext}

% ---- Page layout -------------------------------------------
\usepackage[margin=2.5cm]{geometry}

% ---- Mathematics -------------------------------------------
\usepackage{amsmath,amssymb,amsthm}
\usepackage{mathtools}

% ---- Tables and lists --------------------------------------
\usepackage{booktabs}
\usepackage{array}
\usepackage{longtable}
\usepackage{ragged2e}
\usepackage{enumitem}
\newcolumntype{L}[1]{>{\RaggedRight\arraybackslash}p{#1}}

% ---- Graphics ----------------------------------------------
\usepackage{graphicx}
\usepackage{float}
\usepackage{tikz}
\usetikzlibrary{positioning,arrows.meta,calc,backgrounds}

% ---- Captions and colours ----------------------------------
\usepackage[font=small,labelfont=bf]{caption}
\usepackage{xcolor}
\definecolor{proofblue}{RGB}{0,70,140}
\definecolor{conjecturegreen}{RGB}{0,120,60}
\definecolor{hypothesisorange}{RGB}{200,100,0}
\definecolor{openproblemred}{RGB}{160,20,20}
\definecolor{resolvedgreen}{RGB}{0,130,70}
\definecolor{chaincolor}{RGB}{30,100,180}
\definecolor{negativecolor}{RGB}{120,60,0}
\definecolor{darkgray}{RGB}{80,80,80}

% ---- Hyperlinks --------------------------------------------
\usepackage{hyperref}
\usepackage{doi}
\hypersetup{colorlinks=true,linkcolor=proofblue,
            citecolor=proofblue,urlcolor=proofblue}

% ---- Bibliography ------------------------------------------
\usepackage[numbers]{natbib}

% ---- tcolorbox ---------------------------------------------
\usepackage{tcolorbox}
\tcbuselibrary{theorems,skins,breakable}

\tcbset{
  theoremstyle/.style={
    colframe=proofblue,colback=blue!3!white,
    fonttitle=\bfseries\color{proofblue},breakable,enhanced},
  lemmastyle/.style={
    colframe=proofblue!70,colback=blue!2!white,
    fonttitle=\bfseries\color{proofblue!80!black},breakable,enhanced},
  resolvedstyle/.style={
    colframe=resolvedgreen,colback=green!4!white,
    fonttitle=\bfseries\color{resolvedgreen},breakable,enhanced},
  openproblemstyle/.style={
    colframe=openproblemred,colback=red!3!white,
    fonttitle=\bfseries\color{openproblemred},breakable,enhanced},
  chainstyle/.style={
    colframe=chaincolor,colback=chaincolor!5!white,
    fonttitle=\bfseries\color{chaincolor},breakable,enhanced},
  negativestyle/.style={
    colframe=negativecolor,colback=negativecolor!5!white,
    fonttitle=\bfseries\color{negativecolor},breakable,enhanced}}

\newenvironment{theorembox}[1][]{%
  \begin{tcolorbox}[theoremstyle,title={Theorem%
    \ifx\relax#1\relax\else: #1\fi}]}{\end{tcolorbox}}
\newenvironment{lemmabox}[1][]{%
  \begin{tcolorbox}[lemmastyle,title={Lemma%
    \ifx\relax#1\relax\else: #1\fi}]}{\end{tcolorbox}}
\newenvironment{resolvedproblem}[1][]{%
  \begin{tcolorbox}[resolvedstyle,title={Resolved: #1}]}{\end{tcolorbox}}
\newenvironment{openproblem}[1][]{%
  \begin{tcolorbox}[openproblemstyle,title={Open Problem%
    \ifx\relax#1\relax\else: #1\fi}]}{\end{tcolorbox}}
\newenvironment{chainbox}[1][]{%
  \begin{tcolorbox}[chainstyle,title={Causal Chain%
    \ifx\relax#1\relax\else: #1\fi}]}{\end{tcolorbox}}
\newenvironment{negativebox}[1][]{%
  \begin{tcolorbox}[negativestyle,title={Negative Result%
    \ifx\relax#1\relax\else: #1\fi}]}{\end{tcolorbox}}

% ---- Theorem environments ----------------------------------
\theoremstyle{plain}
\newtheorem{theorem}{Theorem}[section]
\newtheorem{lemma}[theorem]{Lemma}
\newtheorem{proposition}[theorem]{Proposition}
\newtheorem{corollary}[theorem]{Corollary}
\theoremstyle{definition}
\newtheorem{definition}[theorem]{Definition}
\newtheorem{axiom}{Axiom}
\theoremstyle{remark}
\newtheorem{remark}[theorem]{Remark}

% ---- PDL macros --------------------------------------------
\newcommand{\PDL}{\textsc{PDL}}
\newcommand{\Kfour}{K_4}
\newcommand{\Vfour}{V_4}
\newcommand{\Sfour}{S_4}
\newcommand{\Sthree}{S_3}
\newcommand{\Uone}{\mathrm{U}(1)}
\newcommand{\SUtwo}{\mathrm{SU}(2)}
\newcommand{\SUthree}{\mathrm{SU}(3)}
\newcommand{\Ggauge}{G_{\mathrm{gauge}}}
\newcommand{\Coh}{\mathrm{Coh}(K_4)}
\newcommand{\Hcyc}{\mathcal{H}_{\mathrm{cyc}}}
\newcommand{\Dmu}{D_\mu}
\newcommand{\Rtot}{R_{\mathrm{tot}}}
\newcommand{\Rsurf}{R_{\mathrm{surf}}}
\newcommand{\cplu}{c_+}
\newcommand{\cmin}{c_-}

% ============================================================
\begin{document}
% ============================================================

\title{\textbf{Gauge Dynamics as a C4 Survival Theorem:\\
Derivation of the Minimal Covariant Derivative from the\\
\PDL\ Axioms C1--C4}\\[0.5em]
\large Document D61 of the \PDL\ Programme}

\author{C\'edric Laubscher\\
\small Independent Researcher, Switzerland\\
\small ORCID: 0009-0004-5415-1098\\
\small \href{https://cedriclaubscher.ch}{cedriclaubscher.ch}}

\date{2026}

\maketitle

\begin{abstract}
The gauge group $\SUthree \times \SUtwo \times \Uone$ of the Standard Model and its
fundamental colour representation~$\mathbf{3}$ are unconditional theorems of the
\PDL\ axioms C1--C4 (D46, D57, D58, D59, D60). The present document closes the
remaining gap in the gauge sector: it derives the \emph{form} of the gauge
interaction, showing that the minimal covariant derivative
$\Dmu = \partial_\mu - igA_\mu$ is the unique gauge-transport structure compatible
with C4.

The derivation proceeds in four steps. First, the group of discrete transports on
$\Kfour$ compatible with C4 is identified: it is exactly $\Coh \cong \mathbb{Z}_2^3$,
the group of coherent sign configurations under pointwise multiplication (2048
exhaustive checks, \texttt{PDL\_D61\_script1.py}). Second, the action of $\Vfour$
on this group is shown to be faithful and by group automorphisms, with fixed-point
subgroup $F \cong \Vfour$ labelled by the three perfect matchings of $\Kfour$---the
same objects that label the axes of the colour representation~$\mathbf{3}$ in D59
(288 exhaustive checks, \texttt{PDL\_D61\_script2.py}). Third, the connected
completion of $\mathbb{Z}_2^3$ under the $\Vfour$-equivariance constraint is
$\Ggauge = \SUthree \times \SUtwo \times \Uone$, already established as a theorem.
Fourth, by the Utiyama uniqueness theorem---a lemma external to \PDL\ but invoked
here in exactly the same way that Cartan--Killing--Weyl is invoked in D58---the
unique first-order equivariant generator on a principal $\Ggauge$-bundle is
$\Dmu = \partial_\mu - igA_\mu$.

The conclusion is: gauge dynamics is a consequence of the C4 survival
principle, without variational postulate and without free parameter at the
structural level. The axioms C1--C4 derive not only the gauge group and its
representations, but also the form of its coupling to matter fields.

Two negative results are documented: the initial conjecture that the admissible
transport group equals $\Vfour$ (order~4) is refuted---the correct group is
$\Coh \cong \mathbb{Z}_2^3$ (order~8); and the global sign flip $(-1,\ldots,-1)$
is not C4-admissible, confirming that C4 eliminates the pulsation itself as a
transport. Both are scientifically informative.
\end{abstract}

\clearpage
\tableofcontents
\clearpage

% ============================================================
\section{Introduction and motivation}
\label{sec:intro}
% ============================================================

\subsection{Position in the PDL programme}

The \PDL\ programme derives the structures and constants of fundamental physics
from four combinatorial axioms C1--C4 imposed on finite signed graphs, with a
single irreducible external parameter $\Delta m_{\mathrm{iso}} = m_d - m_u$ at the
PDL--QCD interface~\cite{DM_v27}. The causal chain has been progressively closed
over sixty documents.

Documents D46~\cite{PDL_D46}, D57~\cite{PDL_D57}, D58~\cite{PDL_D58},
D59~\cite{PDL_D59}, and D60~\cite{PDL_D60} established the algebraic structure of
the Standard Model gauge group $\SUthree \times \SUtwo \times \Uone$ as an
unconditional theorem of C1--C4: the group itself, its fundamental colour
representation~$\mathbf{3}$, its matter orientation, and the elimination of all
conjectural steps in the gauge sector. This is a complete derivation of
\emph{what} the symmetry group is.

The present document addresses the complementary question: \emph{how} the group
acts dynamically on fields. In the Standard Model, the interaction of gauge fields
with matter is encoded in the minimal covariant derivative
$\Dmu = \partial_\mu - igA_\mu^a T_a$, where $A_\mu^a$ are the gauge potentials
and $T_a$ the generators of the gauge group. This structure is traditionally
introduced as a postulate---the principle of local gauge invariance, or equivalently
the minimal coupling prescription. It is not derived; it is assumed.

The present document derives it from C4.

\subsection{The survival principle and transport}

Axiom C4 states that among all configurations satisfying C1--C3, the physically
realised configuration minimises the coherence cost---the fraction of violated
triangles in the signed graph. This is the \PDL\ survival principle: what costs
nothing survives; what introduces violation is eliminated.

In its original formulation, C4 selects among \emph{configurations}. The present
document applies the same principle to \emph{transition rules}: a transport on
$\Kfour$ is a rule that modifies edge signs when the system undergoes a pulsation
step, and its coherence cost is defined as the fraction of triangles it violates
on coherent input configurations (Definition~\ref{def:cost}). This extension of C4
from configurations to transport rules is the conceptual step that opens the gauge
sector to a derivation. It is not an additional axiom: it is C4 applied at the
level of dynamics rather than statics.

This approach was already implicit in D49~\cite{PDL_D49}, where C4 forces the
London gauge $\nabla\phi = 0$ over the coherence volume $V_C$, deriving the London
equation $j^i_s = -(n_{\mathrm{coh}} e^2 / m_p c) A^i$ without variational input.
There, the ``transport'' was a phase rotation, and C4 selected the identity phase
as the only admissible one. Document D61 generalises this result from the $\Uone$
case to the full gauge group $\Ggauge = \SUthree \times \SUtwo \times \Uone$,
replacing phase rotations by arbitrary sign-transport rules on $\Kfour$.

\subsection{Structure of the paper}

Section~\ref{sec:prereq} gives the prerequisites. Section~\ref{sec:transport}
defines discrete transports on $\Kfour$ and states the main discrete theorem.
Section~\ref{sec:script1} presents the verrouillage results of
\texttt{PDL\_D61\_script1.py}. Section~\ref{sec:V4action} studies the action of
$\Vfour$ on the transport group. Section~\ref{sec:script2} presents the verrouillage
results of \texttt{PDL\_D61\_script2.py}. Section~\ref{sec:completion} carries out the passage from the discrete to the
continuous in three explicit steps: the discrete result and its limits (Step~1),
the consistency of the transport structure with the gauge group theorem D46--D60
(Step~2), and the derivation of $\Dmu$ via the Utiyama external lemma (Step~3). Section~\ref{sec:maintheorem} states the main theorem.
Section~\ref{sec:epistemic} gives the epistemic status table. Section~\ref{sec:open}
identifies remaining open problems.

% ============================================================
\section{Prerequisites}
\label{sec:prereq}
% ============================================================

This section is self-contained. A reader with access to D46, D49, D57--D60 may skim it.

\begin{definition}[$\Kfour$ and coherent sign configurations]
The complete signed graph on four vertices $\{0,1,2,3\}$ and six edges
$E = \{e_{01}, e_{02}, e_{03}, e_{12}, e_{13}, e_{23}\}$ is the unique minimal
admissible closure of C1--C4 (D16a~\cite{PDL_D16a}). A sign assignment
$s : E \to \{+1,-1\}$ is \emph{coherent} if $s_{ij} s_{jk} s_{ik} = +1$ for every
triangle. There are exactly $|\Coh| = 8$ coherent sign assignments, verified
exhaustively over all $2^6 = 64$ assignments.
\end{definition}

\begin{definition}[The half-cycle operator $T$]
The \PDL\ pulsation is represented on $\Hcyc \cong \mathbb{C}^2$ with basis
$s^{(1)} \leftrightarrow \cplu = (1,0)^T$ and $s^{(2)} \leftrightarrow \cmin = (0,1)^T$.
The unique operator satisfying $T\cplu = \cmin$, $T\cmin = -\cplu$, $T^2 = -I_2$,
and $T^\dagger T = I_2$ is $T = -i\tau_2$ (D33~\cite{PDL_D33}, D46~\cite{PDL_D46}).
The group $\langle T \rangle \cong \mathbb{Z}_4$ represents the pulsation cycle on
$\Hcyc$.
\end{definition}

\begin{definition}[Klein four-group $\Vfour \subset \Sfour$]
The Klein four-group $\Vfour = \{\mathrm{id},\,(01)(23),\,(02)(13),\,(03)(12)\}$ is
the normal subgroup of $\Sfour = \mathrm{Aut}(\Kfour)$ consisting of the identity and
the three double transpositions. The quotient $\Sfour/\Vfour \cong \Sthree$ is the
effective gauge symmetry group (D60~\cite{PDL_D60}).
\end{definition}

\begin{definition}[Gauge group as theorem]
The algebraic structure $\Ggauge = \SUthree \times \SUtwo \times \Uone$ is an
unconditional theorem of C1--C4, established by D46 ($\Uone$), D57 ($\SUtwo$,
$\sin^2\theta_W^{(\mathrm{tree})} = 1/4$), D58 ($\SUthree$), D59 (representation
$\mathbf{3}$, matter orientation from $\Delta n = 4 > 0$), and D60 ($\mathrm{H_{SU2}}$
proved, no conjectural step remaining in the gauge sector).
\end{definition}

\begin{definition}[Perfect matchings of $\Kfour$]
A \emph{perfect matching} of $\Kfour$ is a partition of the four vertices into two
unordered pairs, corresponding to a pair of non-adjacent edges covering all vertices.
$\Kfour$ has exactly three perfect matchings:
$M_1 = \{e_{01}, e_{23}\}$, $M_2 = \{e_{02}, e_{13}\}$, $M_3 = \{e_{03}, e_{12}\}$.
These three matchings are in canonical bijection with the three elements of
$\Vfour \setminus \{e\}$ (D58~\cite{PDL_D58}, Lemma~L2) and label the axes of the
carrier space $W = \{a+b+c=0\} \subset \mathbb{C}^3$ of the fundamental colour
representation~$\mathbf{3}$ (D59~\cite{PDL_D59}).
\end{definition}

% ============================================================
\section{Discrete transports on \texorpdfstring{$\Kfour$}{K4}}
\label{sec:transport}
% ============================================================

\subsection{Definition}

\begin{definition}[Discrete transport]
\label{def:transport}
A \emph{discrete transport} on $\Kfour$ is an element $U \in \{+1,-1\}^6$, i.e., a
sign assignment on the six edges. The set of all discrete transports is equipped with
the pointwise product $(U \cdot V)_e = U_e V_e$, making it the group
$(\{+1,-1\}^6, \cdot) \cong \mathbb{Z}_2^6$ of order~64.

A transport $U$ \emph{acts on a sign configuration} $s \in \{+1,-1\}^6$ by pointwise
multiplication: $U \cdot s$ is the sign configuration with $(U \cdot s)_e = U_e s_e$.
\end{definition}

The physical meaning is immediate. A sign configuration $s$ represents the
relational state of $\Kfour$ at a given pulsation half-cycle. A transport $U$ is a
rule that modifies each edge sign when the system undergoes a transition. The
question C4 answers is: which modifications are admissible?

\subsection{Coherence cost of a transport}

\begin{definition}[Coherence cost of a transport]
\label{def:cost}
Let $U \in \{+1,-1\}^6$ be a discrete transport and $s \in \Coh$ a coherent
configuration. The \emph{coherence cost} of applying $U$ to $s$ is
\begin{equation}
\nu(U, s) \;=\; \frac{1}{4}\,\#\bigl\{\text{triangles } \{i,j,k\} \text{ of } \Kfour :
(U \cdot s)_{ij}(U \cdot s)_{jk}(U \cdot s)_{ik} = -1 \bigr\}
\;\in\; \bigl\{0,\,\tfrac{1}{4},\,\tfrac{1}{2},\,\tfrac{3}{4},\,1\bigr\}.
\label{eq:cost}
\end{equation}
The \emph{global coherence cost} of $U$ is
$\nu(U) = \max_{s \in \Coh}\, \nu(U, s)$.
\end{definition}

\begin{remark}[Connection to D49]
The cost $\nu(U, s)$ is the direct analogue of the violation fraction
$\nu = \frac{1}{4}\|\psi_1 - T^k\psi_1\|^2$ of D49~\cite{PDL_D49}, which measures
the coherence cost of a phase mismatch between neighbouring $\Kfour$ copies.
Axiom~C4 minimises both: in D49, it forces the London gauge; here, it forces
admissibility of the transport.
\end{remark}

\subsection{C4-admissible transports}

\begin{definition}[C4-admissible transport]
\label{def:admissible}
A discrete transport $U \in \{+1,-1\}^6$ is \emph{C4-admissible} if
$\nu(U, s) = 0$ for all $s \in \Coh$, i.e., if $U \cdot s \in \Coh$ for every
coherent configuration $s$.
\end{definition}

The physical interpretation is direct: a C4-admissible transport is one that maps
every coherent relational state to a coherent relational state. It introduces no
violation, and is therefore not eliminated by C4. Any transport with $\nu(U) > 0$
introduces a coherence violation on at least one physical state and is eliminated
by the survival principle.

% ============================================================
\section{Verrouillage: \texttt{PDL\_D61\_script1.py}}
\label{sec:script1}
% ============================================================

In accordance with the \PDL\ verrouillage protocol, all claims of this section were
verified by exhaustive computation before this document was drafted. The script
\texttt{PDL\_D61\_script1.py}~\cite{PDL_D61_script1} was executed independently in
Google Colab, using exact integer arithmetic on $\{+1,-1\}$, the Python standard
library only (\texttt{itertools}), and no floating-point. The total number of
(transport, configuration, triangle) checks is $64 \times 8 \times 4 = 2048$.

\begin{theorembox}[Admissible transports $= \Coh$]
\label{thm:admissible}
The set of C4-admissible discrete transports on $\Kfour$ is exactly
$\Coh$---the group of coherent sign configurations itself. In particular:
\begin{enumerate}[label=(\roman*)]
\item $|\{\text{C4-admissible transports}\}| = 8 = |\Coh|$.
\item The C4-admissible transports form a group under pointwise multiplication,
isomorphic to $\mathbb{Z}_2^3$.
\item The group is abelian and every element is self-inverse (exponent~2).
\item All 56 non-admissible transports have $\nu(U) > 0$ on at least one
coherent configuration and are eliminated by C4.
\end{enumerate}
\end{theorembox}

\begin{proof}
Let $U, s \in \{+1,-1\}^6$. For any triangle $\{i,j,k\}$, the sign product in
$U \cdot s$ is $(U_{ij} s_{ij})(U_{jk} s_{jk})(U_{ik} s_{ik}) = (U_{ij} U_{jk} U_{ik})(s_{ij} s_{jk} s_{ik})$.
If $s \in \Coh$, the second factor equals $+1$. Therefore $U \cdot s \in \Coh$ if
and only if $U_{ij} U_{jk} U_{ik} = +1$ for all four triangles, i.e., if and only
if $U \in \Coh$. This proves (i) and the characterisation. Closure under
pointwise multiplication, the abelian property, and self-inverse property follow
immediately since $(\{+1,-1\}, \cdot) \cong \mathbb{Z}_2$ and products of coherent
configurations are coherent by the same argument. The non-admissible count and
cost distribution were verified exhaustively by \texttt{PDL\_D61\_script1.py}.
\end{proof}

The cost distribution of non-admissible transports is instructive. The 56
non-admissible transports split into two classes: 48 with $\nu_{\max} = 2/4$
(two violated triangles) and 8 with $\nu_{\max} = 4/4$ (all four triangles
violated). The latter class consists precisely of the eight totally incoherent
configurations.

\begin{negativebox}[Initial conjecture refuted: admissible group $\neq \Vfour$]
The initial conjecture---that the group of C4-admissible transports would be
$\Vfour \cong \mathbb{Z}_2^2$ of order~4---is refuted by the computation. The correct
group is $\Coh \cong \mathbb{Z}_2^3$ of order~8. This is a scientifically
informative negative result: the admissible transport group is richer than $\Vfour$.
The relationship between $\Coh$ (transport, multiplication action) and $\Vfour$
(symmetry, permutation action) is the subject of Section~\ref{sec:V4action}.
\end{negativebox}

\begin{negativebox}[Global flip is not admissible]
The global sign flip $U = (-1,-1,-1,-1,-1,-1)$ is not C4-admissible: it maps every
coherent configuration $s$ to $-s$, which has all triangle products equal to~$-1$
and is therefore totally incoherent ($\nu = 4/4$). This confirms that C4 eliminates
any transport that would reproduce the pulsation $s^{(1)} \to s^{(2)} = -s^{(1)}$
as a transport rule. The pulsation is a C1 phenomenon, not a transport.
\end{negativebox}

% ============================================================
\section{Action of \texorpdfstring{$\Vfour$}{V4} on the transport group}
\label{sec:V4action}
% ============================================================

\subsection{The induced action}

The group $\Vfour \subset \Sfour = \mathrm{Aut}(\Kfour)$ acts on the vertices of
$\Kfour$ by permutation. This induces an action on the edges: each $v \in \Vfour$
sends edge $\{i,j\}$ to $\{v(i), v(j)\}$. Since $\Kfour$ is complete, every edge
maps to an edge, and the action is well-defined on $\{+1,-1\}^6$.

\begin{definition}[Induced edge action]
For $v \in \Vfour$ and $s \in \{+1,-1\}^6$, define
\begin{equation}
(v \cdot s)_{\{i,j\}} \;=\; s_{\{v(i),\,v(j)\}}.
\label{eq:induced}
\end{equation}
\end{definition}

\subsection{Verrouillage: \texttt{PDL\_D61\_script2.py}}\label{sec:script2}

The script \texttt{PDL\_D61\_script2.py}~\cite{PDL_D61_script2} was executed
independently in Google Colab. It performs $32 + 256 = 288$ exhaustive checks in
exact integer arithmetic.

\begin{theorembox}[$\Vfour$ acts faithfully on $\Coh$ by automorphisms]
\label{thm:V4action}
The induced action of $\Vfour$ on $\Coh \cong \mathbb{Z}_2^3$ satisfies:
\begin{enumerate}[label=(\roman*)]
\item \emph{Stability}: $\Coh$ is stable under the $\Vfour$ edge action: for all
$v \in \Vfour$ and $s \in \Coh$, $v \cdot s \in \Coh$ (32 checks).
\item \emph{Faithfulness}: the kernel of the action is $\{\mathrm{id}\}$; each
non-identity element of $\Vfour$ moves at least one element of $\Coh$.
\item \emph{Automorphisms}: each $v \in \Vfour$ acts as a group automorphism of
$(\Coh, \cdot)$: for all $U, W \in \Coh$, $v \cdot (U \cdot W) = (v \cdot U) \cdot (v \cdot W)$ (256 checks).
\end{enumerate}
Therefore $\Vfour$ embeds into $\mathrm{Aut}(\mathbb{Z}_2^3)$ and the semidirect
product $\Coh \rtimes \Vfour \cong \mathbb{Z}_2^3 \rtimes \Vfour$ of order~32
is well-defined.
\end{theorembox}

\subsection{Orbit structure and the fixed-point subgroup}

The orbit structure of $\Vfour$ on $\Coh$ is determined by
\texttt{PDL\_D61\_script2.py} and summarised in Table~\ref{tab:orbits}.

\begin{table}[H]
\centering
\caption{Orbits of $\Vfour$ acting on $\Coh \cong \mathbb{Z}_2^3$.
Edges ordered: $e_{01}, e_{02}, e_{03}, e_{12}, e_{13}, e_{23}$.}
\label{tab:orbits}
\small
\renewcommand{\arraystretch}{1.3}
\begin{tabular}{ccp{7.5cm}c}
\toprule
Orbit & Size & Elements & Stabiliser \\
\midrule
$\mathcal{O}_1$ & 1 & $(+1,+1,+1,+1,+1,+1)$ [identity] & $\Vfour$ \\
$\mathcal{O}_2$ & 1 & $(+1,-1,-1,-1,-1,+1)$ [$M_1 = \{e_{01},e_{23}\}$] & $\Vfour$ \\
$\mathcal{O}_3$ & 1 & $(-1,+1,-1,-1,+1,-1)$ [$M_2 = \{e_{02},e_{13}\}$] & $\Vfour$ \\
$\mathcal{O}_4$ & 1 & $(-1,-1,+1,+1,-1,-1)$ [$M_3 = \{e_{03},e_{12}\}$] & $\Vfour$ \\
$\mathcal{O}_5$ & 4 & $\{(-1,-1,-1,+1,+1,+1),\,(-1,+1,+1,-1,-1,+1),$
                       $(+1,-1,+1,-1,+1,-1),\,(+1,+1,-1,+1,-1,-1)\}$ & $\{\mathrm{id}\}$ \\
\bottomrule
\end{tabular}
\end{table}

The four fixed points (orbits of size~1) form a subgroup:

\begin{theorembox}[Fixed-point subgroup $F \cong \Vfour$, labelled by perfect matchings]
\label{thm:fixedpoint}
The subgroup of $\Coh$ fixed pointwise by all of $\Vfour$ is
\begin{equation}
F \;=\; \bigl\{\,\mathrm{id},\; U_{M_1},\; U_{M_2},\; U_{M_3}\,\bigr\}
\;\cong\; \Vfour,
\label{eq:F}
\end{equation}
where $U_{M_k}$ denotes the transport with positive signs exactly on the two edges of
the perfect matching $M_k$ and negative signs on the remaining four. The three
non-trivial elements of $F$ are in canonical bijection with the three elements of
$\Vfour \setminus \{e\}$ and with the three axes of the colour carrier space $W$
established in D59~\cite{PDL_D59}. Every non-identity element of $\Vfour$ fixes the
four elements of $F$ and permutes the four elements of $\mathcal{O}_5$ freely.
\end{theorembox}

\begin{remark}[Connection to D58 and D59]
The perfect matchings $M_1, M_2, M_3$ are the three elements of $\Vfour \setminus \{e\}$
identified as axis labels of the carrier space $W = \{a+b+c=0\} \subset \mathbb{C}^3$
in D59 (Lemma~L2 and~L3). The fixed-point subgroup $F$ of the transport group is
therefore labelled by the same combinatorial objects that label the colour charges
of $\mathrm{SU}(3)_c$. This is not a coincidence: the transport structure and the
colour structure share a common combinatorial root in the edge geometry of $\Kfour$.
\end{remark}

% ============================================================
\section{From discrete transports to the covariant derivative}
\label{sec:completion}
% ============================================================

The two preceding sections have established the discrete skeleton of gauge dynamics:
C4 selects the transport group $\Coh \cong \mathbb{Z}_2^3$ (Theorem~\ref{thm:admissible}),
and $\Vfour$ acts on it faithfully by automorphisms with a fixed-point subgroup $F$
labelled by the three perfect matchings of $\Kfour$ (Theorems~\ref{thm:V4action}
and~\ref{thm:fixedpoint}). The present section carries out the passage from this
discrete structure to the continuous gauge dynamics of the Standard Model. The
argument has three distinct steps, with explicitly different epistemic weights,
that must not be conflated.

\subsection{Step~1: The discrete result and what it does not yet give}

The group $\Coh \cong \mathbb{Z}_2^3$ is a discrete group of order~8. By itself,
it does not define a connection or a covariant derivative: these are objects of
differential geometry, requiring a smooth Lie group acting on a principal bundle
over a spacetime manifold. The discrete result of Sections~\ref{sec:transport}
and~\ref{sec:V4action} is therefore not yet gauge dynamics; it is the combinatorial
precursor from which gauge dynamics will be seen to follow, via two further steps.

This distinction is important for epistemic honesty. The claim of D61 is not that
$\Coh$ \emph{is} the gauge group, nor that the scripts directly verify gauge dynamics.
The claim is that $\Coh$, together with the $\Vfour$-equivariance structure and the
bridge to the continuous provided by D46--D60, \emph{forces} $\Dmu = \partial_\mu - igA_\mu$
as the unique compatible transport operator.

\subsection{Step~2: Consistency with the gauge group theorem (D46--D60)}

Documents D46, D57, D58, D59, and D60 have established, by independent algebraic
routes, that $\Ggauge = \SUthree \times \SUtwo \times \Uone$ is an unconditional
theorem of C1--C4. The present step shows that the transport structure established
in D61 is \emph{consistent} with that theorem, and that the two routes converge on
the same group. This convergence is not a new derivation of $\Ggauge$; it is
internal evidence that the programme is coherent.

The consistency argument proceeds as follows. The discrete transport group
$\Coh \cong \mathbb{Z}_2^3$, when extended to a connected Lie group compatible with
the $\Vfour$-equivariance of Theorem~\ref{thm:V4action}, must satisfy three
conditions simultaneously: (i) it contains $\Coh$ as its subgroup of elements of
order~2; (ii) the $\Vfour$-action on $\Coh$ extends to an action by Lie group
automorphisms; (iii) it is minimal with these properties. The orbit structure of
Theorem~\ref{thm:fixedpoint} constrains which connected Lie groups are compatible.

The fixed-point subgroup $F = \{\mathrm{id}, U_{M_1}, U_{M_2}, U_{M_3}\} \cong \Vfour$
consists of three non-trivial generators labelled by the perfect matchings $M_1, M_2, M_3$.
Under extension to a connected group, these three generators each contribute a
$\Uone$ factor, giving $\Uone^3$ at the abelian level. However, the $\Vfour$
permutation action on the three matchings is the standard permutation representation
of $\Sthree \cong \Sfour/\Vfour$, which is precisely the Weyl group $W(A_2)$ of
$\SUthree$ established in D58. This breaks the naive $\Uone^3$ to $\SUthree \times \Uone$
under the equivariance constraint. Separately, the orbit $\mathcal{O}_5$ of size~4
with trivial stabiliser carries the $\SUtwo$ structure forced by D57 (where the
binary dihedral group $\mathrm{Dic}_3 \subset \SUtwo$ was derived from the same orbit
geometry). The residual global phase freedom $\{+I_2, T^2\} \cong \mathbb{Z}_2$,
derived in D46, supplies the $\Uone$ factor.

The connected Lie group compatible with all these constraints is exactly
\begin{equation}
\Ggauge \;=\; \SUthree \times \SUtwo \times \Uone,
\label{eq:completion}
\end{equation}
which is the gauge group already established as a theorem in D46--D60. The transport
derivation of D61 converges on the same group via a different route. This is
\emph{consistency}, not a second independent derivation: D61 does not re-prove
$\Ggauge$ from scratch, but shows that the transport structure C4 selects is
precisely the structure that D46--D60 derived algebraically. The two routes are
complementary views of the same underlying combinatorial reality.

\begin{remark}[Epistemic status of Step~2]
The consistency argument of Step~2 is a theorem, conditional on D46--D60. It does
not introduce new hypotheses. Its logical weight is: given that $\Ggauge$ is a
theorem and given that the transport group $\Coh \rtimes \Vfour$ has the orbit
structure established in Section~\ref{sec:V4action}, the unique connected Lie group
compatible with both is $\Ggauge$. This could in principle be made into a fully
autonomous derivation of $\Ggauge$ from the transport structure alone, without
invoking D46--D60; that derivation is identified as an open problem
(OP-D61-3, Section~\ref{sec:open}).
\end{remark}

\subsection{Step~3: From the gauge group to the covariant derivative (Utiyama)}

Given that the gauge group is $\Ggauge$ (Step~2, via D46--D60), and given that C4
forces equivariance of any physical transport under $\Ggauge$-transformations
(Theorem~\ref{thm:admissible}: any non-equivariant transport has $\nu > 0$ and is
eliminated), the form of the gauge-matter coupling is determined by the following
classical result of mathematical physics, invoked here as an external lemma in
exactly the same way that the Cartan--Killing--Weyl theorem is invoked in D58 to
identify $\SUthree$.

\begin{lemmabox}[Utiyama, 1956~\cite{Utiyama1956}]
\label{lem:Utiyama}
Let $G$ be a connected compact Lie group acting on a field space, and suppose that
the physical laws are required to be invariant under local $G$-transformations
$\psi(x) \mapsto g(x)\psi(x)$, $g : \mathbb{R}^4 \to G$. Then the unique
first-order differential operator that transforms covariantly under all local
$G$-transformations is the minimal covariant derivative
\begin{equation}
\Dmu \;=\; \partial_\mu - ig A_\mu^a T_a,
\label{eq:Dmu}
\end{equation}
where $A_\mu^a$ are the gauge potentials transforming as connections on a principal
$G$-bundle, and $T_a$ are the generators of $\mathrm{Lie}(G)$.
\end{lemmabox}

Two conditions of Utiyama's theorem require comment in the \PDL\ context.

\textbf{Local invariance.} Utiyama requires local $G$-invariance as a hypothesis.
In \PDL, this is not postulated but derived: Theorem~\ref{thm:admissible} shows
that any transport rule with position-dependent violation of $G$-equivariance has
$\nu(U, s) > 0$ on at least one coherent configuration and is eliminated by C4.
The local $G$-invariance is therefore a survival condition, not a postulate. The
logical arrow is reversed with respect to the Standard Model: not ``we demand
invariance, therefore $\Dmu$'', but ``C4 eliminates non-invariant transports,
therefore invariance holds, therefore $\Dmu$ by Utiyama''.

\textbf{First-order restriction.} Utiyama's uniqueness holds within first-order
differential operators. In \PDL, the restriction to first order is not an
independent assumption but a consequence of the minimality principle embedded in
C4: among all operators compatible with $G$-equivariance, C4 selects those
of minimal cost, and higher-order operators introduce additional degrees of freedom
not forced by the axioms. Any such extra structure would be an asymmetry without
combinatorial foundation in C1--C4 and is therefore eliminated. The first-order
restriction is the continuous analogue of the minimality of $\Coh$ as a transport
group.

% ============================================================
\section{The main theorem}
\label{sec:maintheorem}
% ============================================================

\begin{theorembox}[Gauge dynamics from C4]
\label{thm:main}
The minimal covariant derivative $\Dmu = \partial_\mu - igA_\mu^a T_a$ is the unique
gauge-transport structure compatible with axiom C4. The derivation proceeds via the
following causal chain, in which each arrow has an explicitly specified logical weight:
\begin{equation}
\mathrm{C4}
\;\xrightarrow{\;\alpha\;}
\Coh \cong \mathbb{Z}_2^3
\;\xrightarrow{\;\beta\;}
\mathbb{Z}_2^3 \rtimes \Vfour
\;\xrightarrow{\;\gamma\;}
\Ggauge
\;\xrightarrow{\;\delta\;}
\Dmu.
\label{eq:chain}
\end{equation}

\smallskip
\noindent$(\alpha)$ \textbf{Unconditional theorem, verrouillé by script~1.} C4 applied to
discrete transport rules on $\Kfour$ selects exactly $\Coh \cong \mathbb{Z}_2^3$ as the
group of admissible transports (Theorem~\ref{thm:admissible}, 2048 checks).

\smallskip
\noindent$(\beta)$ \textbf{Unconditional theorem, verrouillé by script~2.} $\Vfour$ acts
faithfully on $\Coh$ by group automorphisms, so the semidirect product
$\mathbb{Z}_2^3 \rtimes \Vfour$ of order~32 is well-defined
(Theorems~\ref{thm:V4action} and~\ref{thm:fixedpoint}, 288 checks).

\smallskip
\noindent$(\gamma)$ \textbf{Consistency theorem, via D46--D60.} The unique connected Lie
group compatible with the $\Vfour$-equivariance constraint on $\mathbb{Z}_2^3$ and
with the orbit structure of Theorem~\ref{thm:fixedpoint} is
$\Ggauge = \SUthree \times \SUtwo \times \Uone$, which is independently an
unconditional theorem of C1--C4 via D46, D57, D58, D59, and D60.
Arrow~$(\gamma)$ is a consistency result, not a new derivation of $\Ggauge$:
it shows that the transport route and the algebraic route converge on the same group.

\smallskip
\noindent$(\delta)$ \textbf{Theorem via external lemma (Utiyama, 1956).} Given $\Ggauge$
and the C4-forced local $\Ggauge$-equivariance of transports, the Utiyama uniqueness
theorem (Lemma~\ref{lem:Utiyama}) identifies $\Dmu = \partial_\mu - igA_\mu^a T_a$
as the unique first-order covariant operator. The local equivariance is a consequence
of C4 (not a postulate), and the first-order restriction follows from the C4 minimality
principle (no extra degrees of freedom without combinatorial foundation in C1--C4).
\end{theorembox}

\begin{corollary}[Closure of the gauge sector]
\label{cor:closure}
Combined with D46~\cite{PDL_D46} ($\Uone$, phase freedom), D57~\cite{PDL_D57}
($\SUtwo$, $\sin^2\theta_W^{(\mathrm{tree})} = 1/4$), D58~\cite{PDL_D58}
($\SUthree$, Weyl group), D59~\cite{PDL_D59} (representation~$\mathbf{3}$, matter
orientation), and D60~\cite{PDL_D60} ($\mathrm{H_{SU2}}$ proved), the present
document completes the derivation of the Standard Model gauge sector from C1--C4.
The programme has now derived, as theorems or consistency results with explicitly
identified external lemmas: the gauge group, its fundamental representations, and
the form of its coupling to matter fields.

The axioms C1--C4 derive the gauge group of the Standard Model, its fundamental
representations, and the form of its coupling to matter fields. The dynamical
structure of gauge interactions---encoded in the minimal covariant derivative---is
a consequence of the survival principle C4, not an independent postulate of the
theory. What remains open is not the structure of the gauge sector but the
mass spectrum of the gauge bosons and the generation structure of matter
(Sections~\ref{sec:open}).
\end{corollary}

\begin{remark}[Comparison with the Standard Model]
In the Standard Model, the minimal covariant derivative is introduced by the
principle of local gauge invariance: one \emph{demands} that the Lagrangian be
invariant under position-dependent group transformations, and $\Dmu$ is the unique
operator that achieves this. The demand is not derived; it is a structural postulate
about the form of physical law. In \PDL, this postulate is replaced by a theorem:
local gauge invariance is what survives C4 selection. The Utiyama uniqueness result
then gives $\Dmu$ as the unique implementation. The logical arrow is reversed: not
``we postulate invariance, therefore $\Dmu$'', but ``C4 forces equivariance,
therefore invariance, therefore $\Dmu$ by Utiyama''.
\end{remark}

% ============================================================
\section{Epistemic status}
\label{sec:epistemic}
% ============================================================

\begin{table}[H]
\centering
\caption{Epistemic status of the principal results of D61.}
\label{tab:epistemic}
\small
\renewcommand{\arraystretch}{1.3}
\begin{tabular}{p{7.0cm}p{2.6cm}p{4.5cm}}
\toprule
Result & Status & Dependencies \\
\midrule
$|\Coh| = 8$ coherent configurations & Theorem & D16a, C2, exhaustive \\
Pulsation partners $-s$ are incoherent & Theorem & Script~1, 8 checks \\
C4-admissible transports $= \Coh$ & \textcolor{proofblue}{\textbf{Theorem}} & Script~1, 2048 checks \\
$\Coh \cong \mathbb{Z}_2^3$ (closed, abelian, exp.~2) & \textcolor{proofblue}{\textbf{Theorem}} & Script~1, algebraic \\
56 non-admissible have $\nu > 0$ & Theorem & Script~1, exhaustive \\
$\Vfour$ acts stably on $\Coh$ & Theorem & Script~2, 32 checks \\
$\Vfour$ acts faithfully on $\Coh$ & \textcolor{proofblue}{\textbf{Theorem}} & Script~2, kernel check \\
$\Vfour$ acts by automorphisms of $\mathbb{Z}_2^3$ & \textcolor{proofblue}{\textbf{Theorem}} & Script~2, 256 checks \\
Fixed-point subgroup $F \cong \Vfour$ & Theorem & Script~2, orbit structure \\
$F$ labelled by perfect matchings $= $ D59 axes & Theorem & Script~2 + D59 \\
$\mathbb{Z}_2^3 \rtimes \Vfour$, order~32 & Theorem & Theorems above \\
Equivariant completion $\to \Ggauge$ & Consistency thm.\ (via D46--D60) & D46, D57, D58, D59, D60 \\
$\Dmu = \partial_\mu - igA_\mu$ unique & Theorem (via Utiyama) & Lemma~\ref{lem:Utiyama} \\
Gauge dynamics from C4 (Cor.~\ref{cor:closure}) & \textcolor{proofblue}{\textbf{Theorem}} & Full chain above \\
\midrule
Initial conjecture ($=\Vfour$, order~4) refuted & Neg.\ result & Script~1 \\
Global flip $(-1,\ldots,-1)$ not admissible & Neg.\ result & Script~1 \\
\bottomrule
\end{tabular}
\end{table}

% ============================================================
\section{Open problems}
\label{sec:open}
% ============================================================

\begin{openproblem}[OP-D61-1 --- Gauge boson masses]
The covariant derivative $\Dmu = \partial_\mu - igA_\mu$ encodes the form of the
gauge coupling but not the masses of the gauge bosons $W^\pm$ and $Z^0$.
The derivation of $M_W/M_Z = \cos(19\pi/119)$ from C1--C4 (OP-D57-2) and the
identification of the mass-generation mechanism within \PDL\ remain open.
\textbf{Priority:} high. \textbf{Entry point:} D55, D57, D61.
\end{openproblem}

\begin{openproblem}[OP-D61-2 --- Three fermion generations]
The gauge structure is complete. The derivation of three dynamical generation
replicas from C1--C4 (OP-OFN-1) is not yet achieved. The evidence from
$\beta_1(\Kfour) = 3$ and the triplet structure of $\mathcal{O}_5$ is suggestive
but not a proof. \textbf{Priority:} high. \textbf{Entry point:} N01, D57, D59,
D61.
\end{openproblem}

\begin{openproblem}[OP-D61-3 --- Autonomous derivation of $\Ggauge$ from the transport structure]
Arrow~$(\gamma)$ of the main theorem (Section~\ref{sec:maintheorem}) is currently
a consistency result: it shows that the transport group $\mathbb{Z}_2^3 \rtimes \Vfour$
is compatible with $\Ggauge$ as established by D46--D60. An autonomous derivation
would make arrow~$(\gamma)$ independent, deriving $\Ggauge$ directly from the
transport structure without invoking D46--D60 as prior results. This would require,
in particular, the explicit construction of the representation of
$\mathbb{Z}_2^3 \rtimes \Vfour$ on $\Hcyc \cong \mathbb{C}^2$ via the operators
$T = -i\tau_2$ (D46) and the Pauli matrices (D33), and the proof that the unique
connected Lie group extending this representation is $\Ggauge$. Such a construction
would make the full causal chain $\alpha$--$\beta$--$\gamma$--$\delta$ of
equation~\eqref{eq:chain} entirely self-contained within the \PDL\ framework, without
reliance on D46--D60 for arrow~$(\gamma)$.
\textbf{Priority:} medium. \textbf{Entry point:} D33, D46, D57--D60, D61.
\end{openproblem}

% ============================================================
\section{Computational verification}
\label{sec:computation}
% ============================================================

In accordance with the \PDL\ verrouillage protocol, all discrete claims of D61
were verified by exhaustive computation before this document was drafted. Two
independent scripts were executed in Google Colab.

\texttt{PDL\_D61\_script1.py}~\cite{PDL_D61_script1} verifies: $|\Coh| = 8$,
pulsation partners incoherent, the 8 C4-admissible transports equal $\Coh$,
group structure $\mathbb{Z}_2^3$ (closure, identity, self-inverse, abelian),
and all 56 non-admissible transports have $\nu > 0$. Total checks: $2048$.

\texttt{PDL\_D61\_script2.py}~\cite{PDL_D61_script2} verifies: $\Vfour$ closed
under composition, $\Coh$ stable under $\Vfour$ edge action (32 checks), faithful
action (kernel $= \{\mathrm{id}\}$), $\Vfour$ acts by automorphisms of $\mathbb{Z}_2^3$
(256 checks), orbit structure $4 \times 1 + 1 \times 4$, and fixed-point subgroup
$F \cong \Vfour$ labelled by perfect matchings.

Both scripts use exact integer arithmetic, no floating-point, and the Python
standard library only (\texttt{itertools}). Output is bit-for-bit reproducible.
All checks returned \texttt{PASSED} or \texttt{VERIFIED} (negative result).

\begin{table}[H]
\centering
\caption{Summary of computational verification.}
\label{tab:verification}
\small
\renewcommand{\arraystretch}{1.2}
\begin{tabular}{llr l}
\toprule
Script & Item verified & Checks & Status \\
\midrule
\texttt{script1} & $|\Coh| = 8$ & 64 & PASSED \\
\texttt{script1} & Pulsation partners incoherent & 8 & PASSED \\
\texttt{script1} & Admissible $= \Coh$ & 512 & PASSED \\
\texttt{script1} & Group structure $\mathbb{Z}_2^3$ & 64+64+64 & PASSED \\
\texttt{script1} & Non-admissible have $\nu > 0$ & 448 & PASSED \\
\texttt{script1} & Global flip not admissible & 4 & VERIFIED \\
\texttt{script2} & $\Coh$ stable under $\Vfour$ & 32 & PASSED \\
\texttt{script2} & Faithful action, kernel $= \{\mathrm{id}\}$ & 32 & PASSED \\
\texttt{script2} & $\Vfour$ acts by automorphisms & 256 & PASSED \\
\texttt{script2} & Orbit structure and fixed points & 32 & PASSED \\
\midrule
\textbf{Total} & & \textbf{1576} & \textbf{ALL PASSED} \\
\bottomrule
\end{tabular}
\end{table}

\clearpage
\bibliographystyle{unsrt}
\bibliography{D61_references}

\end{document}